
\chapter*{Abstract}
\addcontentsline{toc}{chapter}{Abstract}

This dissertation addresses the critical challenge of robust Speech Emotion Recognition (SER) in the presence of environmental noise. While deep learning has significantly advanced SER, model performance often degrades when moving from clean, laboratory-recorded data to more realistic, noisy conditions. This study conducts a systematic evaluation of deep learning models and audio features to improve noise robustness and generalization. We investigate four distinct architectures: a Multilayer Perceptron (MLP), a shallow Convolutional Neural Network (CNN), a deeper ResNet-18, and the attention-based Audio Spectrogram Transformer (AST). These models are trained on both traditional Mel-frequency cepstral coefficient (MFCC) features and 2D Mel-spectrograms.

The core of our methodology revolves around a comprehensive on-the-fly data augmentation strategy, where training data is mixed with white noise and natural soundscapes from the ESC-50 dataset at various signal-to-noise ratios (SNRs). We evaluate model performance on two standard benchmarks, RAVDESS and IEMOCAP, under both clean and noisy test conditions. The findings consistently demonstrate that the AST architecture surpasses the convolutional and MLP-based models, particularly when trained with data augmentation. Augmentation not only preserves but often improves AST's performance on clean test sets while significantly enhancing its robustness against a wide range of noises and SNRs.

Furthermore, our analysis reveals a crucial interaction between model architecture and augmentation strategy. While AST effectively leverages diverse acoustic perturbations to learn noise-invariant features, the ResNet-18 model proves more fragile, exhibiting performance degradation under the same conditions. The study also highlights dataset-specific effects, where augmentation leads to a more balanced precision-recall trade-off on the heterogeneous IEMOCAP dataset, particularly for the AST model. Ultimately, this research underscores the synergy between transformer architectures and robust data augmentation, presenting a promising framework for developing more practical and generalizable SER systems.